\clearpage
\SourceLeaf{065}\pdfbookmark[1]{Sectio secunda.}{sectio-secunda}
\begin{center}
{\large SECTIO SECUNDA.}\par\medskip
\emph{Relationes ad locum simplicem in spatio spectantes.}\par\bigskip
\hypertarget{native-065-article047}{}\pdfbookmark[2]{47.}{article-047}47.
\end{center}
In Sectione prima de motu corporum coelestium in orbitis suis actum est, nulla situs, quem hae orbitae in spatio occupant, ratione habita. Ad hunc situm determinandum, quo relationem locorum corporis coelestis ad quaevis alia spatii puncta assignare liceat, manifesto requiritur tum situs plani in quo orbita iacet respectu cuiusdam plani cogniti (e. g. plani orbitae telluris, \emph{eclipticae}), tum situs apsidum in illo plano. Quae quum commodissime ad trigonometriam sphaericam referantur, superficiem sphaericam radio arbitrario circa Solem ut centrum descriptam fingimus, in qua quodvis planum per Solem transiens circulum maximum, quaevis autem recta e Sole ducta punctum depinget. Planis aut rectis per Solem ipsum non transeuntibus plana rectasque parallelas per Solem ducimus, circulosque maximos et puncta in sphaerae superficie his respondentia etiam illa repraesentare concipimus: potest quoque sphaera radio ut vocant infinito magno descripta supponi, in qua plana rectaeque parallelae perinde repraesentantur.

Nisi itaque planum orbitae cum plano eclipticae coincidit, circuli maximi illis planis respondentes (quos etiam simpliciter orbitam et eclipticam vocabimus) duobus punctis se intersecant, quae \emph{nodi} dicuntur; in nodorum altero corpus e Sole visum e regione australi per eclipticam in borealem transibit, in altero ex hac in illam revertet; nodus prior \emph{ascendens}, posterior \emph{descendens} appellatur. Nodorum situs in ecliptica per eorum distantiam ab aequinoctio
\SourceLeaf{066}vernali medio (\emph{longitudinem}) secundum ordinem signorum numeratam assignamus. Sit, in Fig. 1, $\AscendingNode$ nodus ascendens, $A\AscendingNode B$ pars eclipticae, $C\AscendingNode D$ pars orbitae; motus terrae et corporis coelestis fiant in directionibus ab $A$ versus $B$ et a $C$ versus $D$, patetque angulum sphaericum, quem $\AscendingNode D$ facit cum $\AscendingNode B$, a $0$ usque ad $180^0$ crescere posse, neque tamen ultra, quin $\AscendingNode$ nodus ascendens esse desinat: hunc angulum \emph{inclinationem orbitae} ad eclipticam dicimus. Situ plani orbitae per longitudinem nodi atque inclinationem orbitae determinato, nihil aliud iam requiritur, nisi distantia perihelii a nodo ascendente, quam secundum ipsam directionem motus numeramus, adeoque negativam sive inter $180^0$ et $360^0$ assumimus, quoties perihelium ab ecliptica ad austrum situm est. Notentur adhuc expressiones sequentes. Longitudo cuiusvis puncti in circulo orbitae numeratur ab eo puncto, quod retrorsum a nodo ascendente in orbita tantundem distat, quantum aequinoctium vernale ab eodem puncto retrorsum in ecliptica: hinc \emph{longitudo perihelii} erit summa longitudinis nodi et distantiae perihelii a nodo; \emph{longitudo vera} corporis \emph{in orbita} autem summa anomaliae verae et longitudinis perihelii. Denique \emph{longitudo media} vocatur summa anomaliae mediae et longitudinis perihelii: haec postrema expressio manifesto in orbitis ellipticis tantum locum habere potest.
